\section{Introduction and objective}
Assignments~1 and~2 measured sequential dense matrix multiplication; the fastest sequential method found was a
cache-blocked i-k-j kernel with tile $b=96$ \cite{a2}. Its outer loop over rows of $C$ is naturally parallel, because
every row of $C=AB$ depends only on one row of $A$ and on the whole of $B$. This study asks:

\begin{quote}
\emph{On a 4-core / 8-thread laptop CPU, how do runtime, speedup and efficiency of the blocked dense \texttt{float64}
multiplication change with the number of OpenMP threads $p$, and from which matrix size $n$ does parallel execution pay off?}
\end{quote}

The factor changed is the thread count $p\in\{1,2,3,4,6,8,16\}$ (16 exceeds the 8 hardware threads: oversubscription),
crossed with $n\in\{100,250,500,1000,1500,2000,3000,4000\}$. The outcomes are time per multiplication, peak resident memory,
speedup $S(p)=T_{seq}(n)/T_p(n)$ and efficiency $E(p)=S(p)/p$, all computed from \textbf{medians}.
An optional, separate experiment compares scalar, SSE2 and AVX2 code generation of the same source at fixed thread counts.

\paragraph{Hypotheses (fixed before any measurement of the campaign).}
\begin{enumerate}[label=\textbf{H\arabic*}]
\item For $n\ge1000$, $S(p)$ grows almost linearly up to the 4 physical cores, with $E(4)\ge0.7$; it stays below 4 partly
because the single-core turbo frequency is higher than the all-core one.
\item Going from 4 to 8 threads (Hyper-Threading) adds little, between 0 and $30\%$, because the two hardware threads
of a core share its floating-point units; $E(8)<E(4)$.
\item With $p=16$ (oversubscription) there is no gain over $p=8$ and the time is slightly worse.
\item For small $n$ ($\le250$, kernel below about 1~ms) fork/join and scheduling costs dominate: $S(4)<2$, low efficiency, and
there is a size threshold below which parallel execution does not pay off.
\item The parallel version with one thread takes the same time as the sequential baseline within $\pm5\%$ for $n\ge500$.
\item Peak resident memory is $\approx 3\cdot8n^2$ bytes plus a constant and does not depend on $p$.
\item (Optional extension) With one thread, the SIMD builds are faster than the scalar one (AVX2 $>$ SSE2 $>$ scalar);
with many threads the relative SIMD gain shrinks somewhat.
\end{enumerate}

\section{Background}
\textbf{Blocking.} Tiling the three loops so that a $b\times b$ working set of $A$, $B$ and $C$ is reused from cache reduces
main-memory traffic \cite{lam,goto}; Assignment~2 found $b=96$ best among the tested tiles on its machine \cite{a2}.
\textbf{Shared-memory parallelism.} OpenMP creates a team of threads for a \texttt{parallel} region, which ends with an implicit
barrier \cite{openmp}. If a fraction $s$ of the work is serial, Amdahl's argument bounds the speedup by $1/s$ \cite{amdahl};
fork/join, load imbalance and contention for shared resources (L3 cache, memory bandwidth, the floating-point units of a
core shared by its two hardware threads under Hyper-Threading \cite{marr}) reduce efficiency further. The roofline model
relates attainable throughput to arithmetic intensity and memory bandwidth \cite{roofline}; blocked multiplication has high
intensity, so it is expected to be compute-bound rather than bandwidth-bound. \textbf{Turbo frequency.} The i7-6820HQ has
a base clock of 2.7~GHz and a maximum turbo clock of 3.6~GHz \cite{nbc}; a single active core can run faster than four, so the
ideal $S(p)=p$ is not attainable even for perfectly parallel work. \textbf{Vectorisation.} LLVM's loop vectoriser turns the
contiguous inner $j$ loop into packed SIMD instructions \cite{llvmvec}; with FMA contraction disabled, each element of $C$ still
receives the same additions in the same order, so results stay bit-identical \cite{higham}. \textbf{Reporting.} Medians,
IQRs and speedups from a stated baseline follow the recommendations of \cite{hoefler}.

\section{Methodology}
\subsection{Implementations}
All code is C11 in one source file (\texttt{src/c/matmul3.c}); matrices are row-major \texttt{float64} arrays.
\begin{itemize}
\item \textbf{\texttt{seq} (equivalent baseline).} The blocked kernel of Assignment~2 \cite{a2} (tile $b=96$, i-k-j order inside
a tile), generalised only to work on a range of rows $[r_0,r_1)$ of $C$; the baseline calls it with $[0,n)$, which is exactly the
Assignment~2 kernel. It never enters an OpenMP region and creates no thread (checked: 1 OS thread).
\item \textbf{\texttt{omp} (parallel).} One \texttt{\#pragma omp parallel num\_threads(p)} region; thread $t$ of $q$ computes the
contiguous rows $[\lfloor nt/q\rfloor,\lfloor n(t+1)/q\rfloor)$ of $C$ with the same kernel. Each element of $C$ is written by exactly
one thread, $A$ and $B$ are only read, so there are no data races and no locks; the only synchronisation is the implicit final barrier.
Rows are split evenly instead of distributing $96$-row tiles, which would leave threads idle (e.g.\ $n=1000$ has only 11 tile rows for 8 threads).
\item \textbf{\texttt{ijk} (reference only).} The unchanged i-j-k kernel of Assignment~1 \cite{a1}, run sequentially at $n=500,1000$
to place the study relative to the basic algorithm. It does not replace the equivalent baseline.
\end{itemize}
Because the summation order over $k$ for every $C_{ij}$ is identical in all three kernels and for every $p$, and FMA contraction is
disabled, all methods must return \emph{bit-identical} products; this is tested and checked in every benchmark run.

\textbf{Build.} Apple clang 14.0.0 with Homebrew \texttt{libomp} 23.1.5 (Apple clang ships without an OpenMP runtime):
\\\hspace*{1.5em}\texttt{clang -O2 -std=c11 -ffp-contract=off -Xpreprocessor -fopenmp}\\\hspace*{1.5em}\texttt{-I\$(brew --prefix libomp)/include -L\$(brew --prefix libomp)/lib -lomp}.\\
No \texttt{-march=native} and no \texttt{-ffast-math}, as in Assignments~1--2. At \texttt{-O2} clang already vectorises the inner loop with
128-bit SSE2 instructions; this primary build is the ``SSE2'' build of the extension (byte-identical binary).

\textbf{Thread control.} Each run sets \texttt{OMP\_NUM\_THREADS}$=p$ and \texttt{OMP\_DYNAMIC=false}; the runtime settings printed by
\texttt{OMP\_DISPLAY\_ENV} are stored per run (\texttt{OMP\_PROC\_BIND=false}, \texttt{OMP\_WAIT\_POLICY=PASSIVE}, libomp defaults). The
program reports, for the last call, the OpenMP team size, how many threads received rows, how many \emph{distinct OS thread ids} executed
the team (\texttt{pthread\_threadid\_np}) and the number of OS threads of the process (\texttt{task\_threads}). The C code calls no numerical
library; NumPy, used only for validation outside the measured processes, runs with \texttt{OPENBLAS\_NUM\_THREADS=1}.
\textbf{Affinity:} macOS offers no thread-pinning interface (the Mach affinity policy is only a hint) and libomp does not support binding on
macOS, so threads are \emph{not pinned}; the OS scheduler places them. This is a limitation of the platform, recorded in \texttt{environment.json}.

\subsection{Inputs and environment}
The dense inputs are the Assignment~1/2 matrices: uniform $[0,1)$ \texttt{float64} from NumPy's \texttt{default\_rng} seeded with
$[20261012, n, w]$ ($w=0$ for $A$, $1$ for $B$), written as little-endian binary files. \texttt{gen\_inputs.py} regenerates them and verifies each SHA-256
against the digests of Assignment~2 (all 16 files identical). Table~\ref{tab:env} summarises the machine. The laptop ran on mains power with
all other applications closed, under \texttt{caffeinate -i}.

\begin{table}[!htb]\centering\small
\caption{Environment (from \texttt{results/environment.json}, recorded with \texttt{sysctl}, \texttt{sw\_vers}, \texttt{clang --version}).}
\label{tab:env}\begin{tabular}{lp{11cm}}
\toprule
Item & Value \\
\midrule
CPU & Intel Core i7-6820HQ CPU @ 2.70GHz \\
Cores / threads & 4 physical, 8 logical (Hyper-Threading), 1 package \\
Caches & L1d 32 KiB and L2 256 KiB per core, L3 8 MiB shared, line 64 B \\
RAM & 16 GiB \\
Machine / OS & MacBookPro13,3, macOS 12.7.6 (21H1320) \\
Compiler & Apple clang version 14.0.0 (clang-1400.0.29.202), libomp 23.1.5 \\
Python (validation) & 3.12.7, NumPy 2.4.2 \\
\bottomrule
\end{tabular}

\end{table}

\subsection{Correctness checks before and during benchmarking}
\texttt{test\_correctness.py} runs \texttt{seq}, \texttt{ijk} and \texttt{omp} with $p\in\{1,2,3,4,7,8,16\}$ on: $1\times1$, $2\times2$ and
$3\times3$ products with hand-computed answers, zero and identity matrices on both sides, an all-ones matrix, small-integer matrices (exact
in \texttt{float64}), random matrices for $n\in\{5,17,64,95,96,97,100,191,250\}$ (below, at and above the tile size, not multiples of $p$),
$n=3$ with more threads than rows, seven tile sizes, the reported thread counts, and repeated calls into the same output buffer.
Random cases are compared with NumPy \cite{numpy} using the relative tolerance $2\gamma_n$, $\gamma_n=nu/(1-nu)$, $u=2^{-53}$: each element of
$C$ is a length-$n$ dot product of non-negative values, whose forward error is bounded by $\gamma_n$ relative to the exact value for
either summation order \cite{higham}, so two correct orders differ by at most about $2\gamma_n$. Every product is also compared
\emph{bit for bit} with the \texttt{seq} result. The suite is run on the primary binary and on the three vectorisation builds
(1404 = $4\times351$ checks in total, all passed). During the campaign every process dumps its last product, which the driver validates against
NumPy with $2\gamma_n$ and whose SHA-256 is compared with that of the sequential baseline of the same $n$.

\subsection{Measurement protocol and budgets}
\begin{itemize}
\item \textbf{Timer and boundary.} \texttt{clock\_gettime(CLOCK\_MONOTONIC)} around the multiplication call only. Inside: zero-filling
$C$ (required by the i-k-j update) and, for \texttt{omp}, the fork/join of the thread team. Outside: reading the input files, allocating
and page-touching $C$, checksum, dump and validation. One process per configuration.
\item \textbf{Repetitions.} 1 warm-up (recorded in the raw CSV, excluded from the statistics; it also lets the OpenMP runtime create its
thread pool) and 5 measured samples; 3 samples at $n=4000$ to keep the campaign within its time budget. Small products are batched so
that a sequential sample lasts at least about 0.2~s (batch sizes 607, 39 and 5 for $n=100,250,500$); the reported value is time per
multiplication. Statistics: median and interquartile range (IQR, linear-interpolation quartiles). Speedup and efficiency use medians;
an IQR-based range $[Q_1(T_{seq})/Q_3(T_p),\,Q_3(T_{seq})/Q_1(T_p)]$ is reported as a rough uncertainty.
\item \textbf{Order, pauses and temperature.} Within each $n$ the configurations run in a seeded random order (the sequential baseline first,
so that its SHA-256 is known), with a 3~s pause between processes. Before and after each process the driver stores the
\texttt{CPU\_Speed\_Limit} reported by \texttt{pmset -g therm} (100 = no thermal limit). To estimate drift and between-process variability,
the whole thread sweep at $n=2000$ is \textbf{repeated in a second pass at the end} of the campaign.
\item \textbf{Memory.} Peak resident set size of the process from \texttt{getrusage(RUSAGE\_SELF).ru\_maxrss}, which macOS reports in
\emph{bytes}; it includes inputs, output, code and thread stacks. The driver also polls the RSS with \texttt{ps} every 0.5~s to enforce the budget.
\item \textbf{Budgets (declared before the campaign).} At most 120~s per timed sample and 4~GiB resident memory per process; a configuration
whose $3\cdot8n^2$ array storage exceeds the memory budget is not attempted. Violations would be recorded as \texttt{timeout} /
\texttt{memory\_limit}, never with an invented time. The whole campaign was budgeted at 1--1.5~h; a 37~s pilot (disclosed here, not
reported) was used only to estimate its duration.
\end{itemize}

\subsection{Supplementary experiment: small sizes (added after the main campaign)}
The main campaign showed a large speedup even at the smallest planned size ($n=100$), so the size threshold of H4 was not located. As in
Assignment~2, a supplementary experiment was therefore \emph{added after seeing those results}: campaign \texttt{small} measures $n=16,32,64$ with
\texttt{seq} and \texttt{omp} $p\in\{1,2,4,8\}$ under the same protocol (1 warm-up, 5 samples, batches of 147\,965, 18\,496 and 2\,312
multiplications so that a sequential sample lasts about 0.2~s, randomised order, 3~s pauses, per-run validation). It ran separately, about two hours
after the main campaign (22:19--22:20 UTC), with the laptop again idle and on mains power. The inputs come from the same generator and seed; their SHA-256 digests were
added to the manifest (they did not exist in Assignment~2).

\subsection{Optional vectorisation extension}
The same source is compiled three times with the base flags above, changing only SIMD code generation: \textbf{scalar}
(\texttt{-fno-vectorize -fno-slp-vectorize}), \textbf{SSE2} (default \texttt{-O2}, 128-bit, 2 doubles per instruction) and \textbf{AVX2}
(\texttt{-mavx2}, 256-bit, 4 doubles). FMA stays disabled (\texttt{-ffp-contract=off}), so the three builds must give bit-identical
results, which is tested. Evidence of vectorisation comes from the compiler (\texttt{-Rpass=loop-vectorize}) and from the disassembly of
the kernel (\texttt{otool -tv}): Table~\ref{tab:vecev}. The campaign \texttt{vec} measures each build with \texttt{seq} and \texttt{omp}
$p\in\{1,2,4,8\}$ at $n=1000$ and $2000$ under the same protocol, so that SIMD effects (same $p$, different build) are separated from
thread effects (same build, different $p$).
